\section[普遍局部无环性]{普遍局部无环性（1 月 11 日）}

\setcounter{secnumdepth}{3}

\srcpage{95}

\subsection{回顾}

在最近几讲中，我们讨论了：

\begin{itemize}
\item $D_{\et}(\Bun_G)$。
\item Jacobian 判别法，它被用来通过给出光滑卡来证明 $D_{\et}(\Bun_G)$ 的一些基本性质。
\item ULA 层，一个同时服务于上述两者的技术性概念。例如，它们为 $D_{\et}(\Bun_G)$ 给出了一种「容许性」的概念。
\end{itemize}

\subsection{概形上 ULA 层的背景}

设 $f\colon X\to S$ 是 Noether 概形的有限型分离映射，$A\in D^b_c(X_{\et},\Lambda)$，其中 $n\Lambda=0$，$n\in\cO_S^\times$。

\begin{definition}
称 $A$ 是 \emph{$f$-局部无环的}（$f$-locally acyclic），若对一切落在 $\bar s\to S$ 上方的几何点 $\bar x\to X$，以及一切泛化 $\bar t\rightsquigarrow\bar s$，映射
\[
A_{\bar x}=R\Gamma(X_{\bar x},A)\xrightarrow{\ \sim\ }R\Gamma(X_{\bar x}\times_{S_{\bar s}}\bar t,A)
\]
是同构。
\end{definition}

\begin{remark}
这是拟凝聚层 $\cF/X$ 关于 $S$ 平坦的 étale 类比。下述是 Gabber 最近的一个定理：在某些 Noether 假设下，这一概念在基变换下稳定。
\end{remark}

\begin{theorem}[Gabber]
设 $A$ 是 $f$-局部无环的。则对任意基变换
\[
\begin{tikzcd}
X' \arrow[r, "\tilde g"] \arrow[d, "f'"'] & X \arrow[d, "f"] \\
S' \arrow[r, "g"'] & S
\end{tikzcd}
\]
层 $\tilde g^*A$ 是 $f'$-局部无环的。
\end{theorem}

通常更好的做法是显式要求所有基变换之后的无环性（尽管由 Gabber 定理，这在事后是自动的）。

\begin{definition}
称 $A$ 是 \emph{$f$-普遍局部无环的}（$f$-universally locally acyclic，简记 $f$-ULA），若对任意如上基变换，$\tilde g^*A$ 都是 $f'$-局部无环的。
\end{definition}

\begin{example}
若 $f$ 光滑，则 $\Lambda$（更一般地，任意局部常值层）是 $f$-ULA 的。对 $f=\id\colon X\to X$，逆命题成立：$A$ 是 $\id$-ULA 的当且仅当 $A$ 局部常值。
\end{example}

\begin{lemma}
考虑交换三角
\[
\begin{tikzcd}
& Y \arrow[dl, "h"'] \arrow[dr, "g"] & \\
X \arrow[rr, "f"'] & & S
\end{tikzcd}
\]
其中 $h$ 紧合，且 $B$ 在 $Y$ 上是 $g$-ULA 的。则 $A:=Rh_*B$ 是 $f$-ULA 的。
\end{lemma}

\begin{example}
实践中，$Y\to X$ 常常取作 $f\colon X\to S$ 的奇点消解。例如在仿射 Grassmann 簇的情形中就会如此。
\end{example}

综合例 18.5 与引理 18.6，得到：

\begin{corollary}
若 $f\colon X\to S$ 紧合且 $A$ 是 $f$-ULA 的，则 $Rf_*A$ 局部常值。
\end{corollary}

\srcpage{96}

\begin{remark}[扭曲 Verdier 对偶]
若 $A$ 是 $f$-ULA 的，则 Poincaré 对偶有一个 $A$-扭曲的版本。更确切地说，
\[
\mathbb{D}_{X/S}(A)\otimes f^*B\xrightarrow{\ \sim\ }R\mathrm{Hom}(A,Rf^!B),
\]
其中 $\mathbb{D}_{X/S}(A)=R\mathrm{Hom}(A,Rf^!\Lambda)$。
\end{remark}

若 $A=\Lambda$，这说的是：当 $\Lambda$ 是 $f$-ULA 的（例如 $f$ 光滑时），有 $Rf^!\Lambda\otimes f^*B\xrightarrow{\ \sim\ }Rf^!B$。

\begin{example}
若 $A$ 是 $f$-ULA 的，则有 [LZ20] 最近证明的一个 Verdier 对偶性质：
\[
A\xrightarrow{\ \sim\ }\mathbb{D}_{X/S}(\mathbb{D}_{X/S}(A)).
\]
事实上，他们把 ULA 层刻画为某个对称幺半范畴中的可对偶对象，其中的对偶由 Verdier 对偶给出。
\end{example}

\begin{example}
若 $S$ 是几何点，则所有 $A\in D^b_c(X_{\et},\Lambda)$ 都是 $f$-ULA 的。
\end{example}

\subsection{菱形上的 ULA 层}

我们希望为菱形建立这一概念的变体。一个重要之处在于：对于全无界导出范畴 $D_{\et}(X,\Lambda)$，我们有很好的类比，但「可构造性」是一个微妙的概念。

\begin{example}
设 $i\colon\Spa C\hookrightarrow\mathbb B_C$ 是把一个点映入球的包含。则 $i_*\Lambda$ 不可构造，但它应当是 $\Spa C$ 上的 ULA 层。
\end{example}

\begin{proposition}
若 $X$ 是有限上同调维数的空间菱形（在 $X_{\et}$ 上一致），则 $D_{\et}(X,\Lambda)$ 是紧生成的，且紧对象正是可构造复形\footnote{可构造子集是那些属于由拟紧开子集生成的布尔代数的子集。}（在某个可构造分层之后局部常值）。
\end{proposition}

\begin{example}
设 $j\colon\mathbb T_C=\{T:|T|=1\}\hookrightarrow\{T:|T|\le1\}=\mathbb B_C$。则 $j_!\Lambda$ 可构造。
\end{example}

\begin{definition}
设 $f\colon X\to S$ 是局部空间菱形之间的映射（可紧化、局部具有有限 $\dim.\operatorname{tr.g.}$，使得 $Rf_!$ 有定义）。设 $A\in D_{\et}(X,\Lambda)$。若以下两个条件成立，则称 $A$ 是 \emph{$f$-局部无环的}。
\begin{enumerate}[label=(\arabic*)]
\item 对一切落在 $\bar s\to S$ 上方的几何点 $\bar x\to X$，以及一切泛化 $\bar t\rightsquigarrow\bar s$\footnote{对概形，本可以用 $X_{\bar x}\times_{S_{\bar s}}S_{\bar t}$ 代替 $X_{\bar x}\times_{S_{\bar s}}\bar t$。对 adic 空间，$\bar t$ 没有意义，但 $S_{\bar t}$ 有意义。}，
\[
A_{\bar x}=R\Gamma(X_{\bar x},A)\xrightarrow{\ \sim\ }R\Gamma(X_{\bar x}\times_{S_{\bar s}}S_{\bar t},A).
\]
\item 对一切使 $f\circ j\colon U\to S$ 为 qcqs 的 étale 映射 $j\colon U\to X$，层
\[
R(f\circ j)_!(A|_U)\in D_{\et}(S,\Lambda)
\]
可构造（意即：沿任意如命题 18.13 那样的 $S'\to S$ 拉回之后可构造）。
\end{enumerate}
称 $A$ 是 $f$-ULA 的，若任意基变换之后它都是局部无环的。
\end{definition}

\begin{remark}
对概形，条件 (2) 自动满足，全部信息都在 (1) 中。对菱形则几乎相反：(1) 几乎自动，而 (2) 承载了大部分信息。
\end{remark}

\begin{example}
Gabber 定理的类比对菱形不成立。若 $S=\Spa C$ 且 $X$ 在 $S$ 上上同调光滑，则任意可构造的 $A$ 都是 $f$-局部无环的；但在可构造的 $A$ 中，只有局部常值的那些是 $f$-ULA 的。
\end{example}

\srcpage{97}
\begin{center}
\figP{097A}
\end{center}
\begin{center}
\figP{097B}
\end{center}

\begin{remark}
adic 空间与概形的关键区别在于：在 adic 情形下，严格 Hensel 化要容易得多。具体地说，$X_{\bar x}$ 由 $\Spa(C,C^+)$ 表示，其中 $C$ 是完全代数闭域，$C^+\subset C$ 是赋值子环。特别地，$|X_{\bar x}|$ 是一条全序的点链。
\end{remark}

让我们用它来分析局部无环性定义中的条件 (1)。$|X_{\bar x}|$ 与 $|S_{\bar s}|$ 都是全序的点链，而 $|S_{\bar t}|$ 是 $|S_{\bar s}|$ 的子集。这意味着，对某个落在 $\bar t\rightsquigarrow\bar s$ 上方的特殊化 $\bar y\rightsquigarrow\bar x$，有 $X_{\bar x}\times_{S_{\bar s}}S_{\bar t}=X_{\bar y}$。因此该条件等价于
\[
A_{\bar x}\xrightarrow{\ \sim\ }A_{\bar y}=R\Gamma(X_{\bar y},A)\cong R\Gamma(X_{\bar x}\times_{S_{\bar s}}S_{\bar t},A).
\]
先验地，无环性条件只对某些特殊化提出这一要求。但「(1) 普遍成立」等价于「对一切特殊化 $\bar y\rightsquigarrow\bar x$，映射 $A_{\bar x}\xrightarrow{\sim}A_{\bar y}$ 都是同构」。若后一情形发生，则称 $A$ 是\emph{超收敛的}。

\begin{remark}
$|\mathbb B_C|$ 上的超收敛层等价于其极大 Hausdorff 商 $|\mathbb B_C|$（即与 $\mathbb B_C$ 相伴的 Berkovich 空间）上的层。类似地，超收敛 étale 层等价于该 Berkovich 空间上的 étale 层。
\end{remark}

\srcpage{98}

\subsection{性质}

\begin{lemma}
若 $f$ 上同调光滑且 $A$ 局部常值，则 $A$ 是 $f$-ULA 的。
\end{lemma}

\begin{proof}
$A$ 的超收敛性由局部常值性立得。为验证 (2)，只需看出：若 $f$ 是 qcqs 且上同调光滑，则 $Rf_!$ 保持可构造复形。我们有伴随
\[
\begin{tikzcd}
D_{\et}(X,\Lambda) \arrow[r, bend left, "Rf_!"] & D_{\et}(S,\Lambda) \arrow[l, bend left, "Rf^!"]
\end{tikzcd}
\]
我们知道，可构造复形正是这些范畴中的紧对象。
\end{proof}

\begin{lemma}
设 $F\colon\cC\to\cD$ 与 $G\colon\cD\to\cC$ 是紧生成三角范畴之间的伴随（$F$ 为左伴随，$G$ 为右伴随）。则 $F$ 保持紧对象当且仅当 $G$ 与所有直和交换。
\end{lemma}

\begin{proof}
这通过抽象的图表追踪来证明。例如，若 $G$ 保持直和且 $A\in\cC$ 是紧对象，则
\begin{align*}
\Hom_{\cD}\Bigl(F(A),\bigoplus_{i\in I}B_i\Bigr)
&\cong\Hom_{\cC}\Bigl(A,G\Bigl(\bigoplus_{i\in I}B_i\Bigr)\Bigr)\\
&\cong\Hom_{\cC}\Bigl(A,\bigoplus_{i\in I}G(B_i)\Bigr)\\
&\cong\bigoplus_{i\in I}\Hom_{\cC}(A,G(B_i))\\
&\cong\bigoplus_{i\in I}\Hom_{\cD}(F(A),B_i).
\end{align*}
\end{proof}

\begin{proof}[关于上同调光滑的引理的证明（续）]
为完成该引理的证明，只需论证 $Rf^!$ 与任意直和交换。但 $f$ 上同调光滑蕴含 $Rf^!\cong f^*\otimes Rf^!\Lambda$，后者显然与直和交换。
\end{proof}

\begin{lemma}
若 $f=\id\colon X\to X$，则 $A$ 是 $f$-ULA 的当且仅当 $A$ 局部常值且具有完美纤维。
\end{lemma}

\begin{proof}
设 $A$ 是 $f$-ULA 的。条件 (2) 蕴含 $A$ 可构造，条件 (1) 蕴含 $A$ 超收敛；而「可构造加超收敛蕴含局部常值」是一般事实。
\end{proof}

\begin{lemma}
考虑交换三角
\[
\begin{tikzcd}
& Y \arrow[dl, "h"'] \arrow[dr, "g"] & \\
X \arrow[rr, "f"'] & & S
\end{tikzcd}
\]
其中 $h$ 紧合，且 $B$ 在 $Y$ 上是 $g$-ULA 的。则 $A:=Rh_*B$ 是 $f$-ULA 的。
\end{lemma}

\srcpage{99}

\begin{proof}
直接论证，使用紧合基变换。
\end{proof}

\begin{corollary}
若 $f\colon X\to S$ 紧合且 $A$ 是 $f$-ULA 的，则 $Rf_*A$ 局部常值。
\end{corollary}

\begin{lemma}[扭曲 Verdier 对偶]
若 $A$ 是 $f$-ULA 的，则对一切 $B\in D_{\et}(S,\Lambda)$ 有
\[
\mathbb{D}_{X/S}(A)\otimes f^*B\xrightarrow{\ \sim\ }R\mathrm{Hom}(A,Rf^!B),
\]
其中 $\mathbb{D}_{X/S}(A)=R\mathrm{Hom}(A,Rf^!\Lambda)$。
\end{lemma}

注意，上述公式蕴含 ULA 定义中的条件 (2)。事实上，它蕴含 $R\mathrm{Hom}(A,Rf^!-)\colon D_{\et}(S,\Lambda)\to D_{\et}(X,A)$\yizhu{此处原文写作 $D_{\et}(X,A)$，按上下文应为 $D_{\et}(X,\Lambda)$。}{A（原文 p.99）} 与所有直和交换，因为该式左端的函子是左伴随的复合（即 $f^*$ 与 $\overset{L}{\otimes}_\Lambda$），而它们保持直和。于是由上述紧对象引理，$Rf^!$ 的左伴随保持可构造性。该左伴随是 $Rf_!(A\overset{L}{\otimes}_\Lambda-)$；把它作用于 $j_!\Lambda$，其中 $j\colon U\hookrightarrow X$ 是拟紧开包含。

\begin{lemma}[Verdier 双对偶]
若 $A$ 是 $f$-ULA 的，则 $\mathbb{D}_{X/S}(A)$ 是 $f$-ULA 的，且典范映射
\[
A\xrightarrow{\ \sim\ }\mathbb{D}_{X/S}(\mathbb{D}_{X/S}(A))
\]
是同构。
\end{lemma}

\begin{example}
设 $S=\Spa C$，$X=X_0^\diamond$，其中 $X_0/C$ 是某个代数簇。则对任意 $A_0\in D^b_c(X_{0,\et},\Lambda)$，其解析化 $A\in D_{\et}(X,\Lambda)$ 是 ULA 的。

例如，这说明对 $i\colon\Spa C\hookrightarrow\mathbb A^1_C$，$i_*\Lambda$ 是 ULA 的。回忆一下，它并不可构造。
\end{example}

\subsection{Verdier 双对偶的证明}

引理 18.26 有两个证明，它们都用到 2-范畴中「可对偶性」的概念。

\subsubsection{Lu--Zheng 的进路}

固定基 $S$。考虑如下定义的对称幺半 2-范畴 $LZ_S$：
\begin{itemize}
\item 对象是 $(X,A)$，其中 $X\to S$ 如上，$A\in D_{\et}(X,\Lambda)$。
\item 态射 $(X,A)\to(Y,B)$ 是上同调对应，意即一个图表
\[
\begin{tikzcd}
Z \arrow[r, "c_2"] \arrow[d, "c_1"'] & Y \arrow[d] \\
X \arrow[r] & S
\end{tikzcd}
\]
连同映射 $c_1^*A\to Rc_2^!B$。
\item 对称幺半结构是 $(X,A)\otimes(Y,B):=(X\times_S Y,A\boxtimes B)$。
\end{itemize}

结果表明，这定义了一个闭对称幺半结构，其内 $\Hom$ 为
\[
\Hom_{LZ_S}\bigl((X,A),(Y,B)\bigr)=(X\times_S Y,R\mathrm{Hom}(p_1^*A,Rp_2^!B)).
\]

\begin{theorem}
以下条件等价。
\begin{enumerate}[label=(\arabic*)]
\item $A$ 是 $(X\to S)$-ULA 的。
\srcpage{100}
\item $(X,A)$ 在 $LZ_S$ 中可对偶化。
\item $(X,A)\otimes(X,A)^\vee\xrightarrow{\ \sim\ }\Hom_{LZ_S}\bigl((X,A),(X,A)\bigr)$，即
\[
p_1^*\mathbb{D}_{X/S}(A)\overset{L}{\otimes}_\Lambda p_2^*A\xrightarrow{\ \sim\ }R\mathrm{Hom}_\Lambda(p_1^*A,Rp_2^!A),
\]
其中方块为：
\[
\begin{tikzcd}
X\times_S X \arrow[r, "p_1"] \arrow[d, "p_2"'] & X \arrow[d, "f"] \\
X \arrow[r, "f"'] & S
\end{tikzcd}
\]
\end{enumerate}
\end{theorem}

此时，对偶为 $(X,A)^\vee=(X,\mathbb{D}_{X/S}(A))$，因此 $\mathbb{D}_{X/S}(A)$ 是 $(X\to S)$-ULA 的，且 $A\xrightarrow{\sim}\mathbb{D}_{X/S}(\mathbb{D}_{X/S}(A))$。

注意，上述 (3) 是 $A$-扭曲 Poincaré 对偶的一个实例。

\begin{corollary}
(1) 常层 $\Lambda$ 是 $f$-ULA 的当且仅当 $p_1^*\mathbb{D}_{X/S}\to\mathbb{D}_{X\times_S X/X}=Rp_2^!\Lambda$ 是同构\yizhu{此处原文只写 $\mathbb{D}_{X/S}$ 而未标明作用对象，应理解为对偶化复形 $\mathbb{D}_{X/S}(\Lambda)$。}{A（原文 p.100）}。

(2) $f$ 关于 $\Lambda$ 上同调光滑当且仅当 $\mathbb{D}_{X/S}$ 可逆，且 $p_1^*\mathbb{D}_{X/S}\xrightarrow{\ \sim\ }Rp_2^!\Lambda$。
\end{corollary}

\subsubsection{第二个证明}

如下定义一个 2-范畴 $\cC_S$：
\begin{itemize}
\item 对象是如上 $X\to S$。
\item 态射为 $\mathrm{Fun}_{\cC_S}(X,Y)=D_{\et}(X\times_S Y,\Lambda)$。
\item 复合是卷积。对 $X,Y,Z\to S$，
\[
\begin{tikzcd}
& X\times_S Y\times_S Z \arrow[dl, "\pi_{12}"'] \arrow[d, "\pi_{13}"] \arrow[dr, "\pi_{23}"] & \\
X\times_S Y & X\times_S Z & Y\times_S Z
\end{tikzcd}
\]
对 $A\in\mathrm{Fun}_{\cC_S}(X,Y)$ 与 $B\in\mathrm{Fun}_{\cC_S}(Y,Z)$，定义\yizhu{原文公式中下标印作 $\pi_{B!}$，按定义应为 $\pi_{13!}$，即到 $X\times_S Z$ 的投影。}{A（原文 p.100）}：
\[
A\star B:=R\pi_{B!}(\pi_{12}^*A\overset{L}{\otimes}_\Lambda\pi_{23}^*B).
\]
\end{itemize}

紧合基变换蕴含结合性。恒等态射是 $\Delta_{X/S!}\Lambda$。

这给出到 2-范畴 $\cC'_S$ 的映射，其定义为：
\begin{itemize}
\item 对象是如上 $X\to S$。
\item 态射 $X\to Y$ 是函子 $D_{\et}(X,\Lambda)\to D_{\et}(Y,\Lambda)$。
\end{itemize}

注意，这破坏了 $\cC_S$ 中 $X$ 与 $Y$ 之间的对称性。函子 $\cC_S\to\cC'_S$ 以层为核。

回忆，在任意 2-范畴中都有伴随函子的概念：称 $f\colon X\to Y$ 是 $g\colon Y\to X$ 的左伴随，若存在 $\alpha\colon\id_X\to gf$ 与 $\beta\colon fg\to\id_Y$，使得
\[
f\xrightarrow{f\alpha}fgf\xrightarrow{\beta f}f
\qquad\text{与}\qquad
g\xrightarrow{\alpha g}gfg\xrightarrow{g\beta}g
\]
同构于相应的恒等自然变换。

\begin{theorem}
以下条件等价：
\begin{enumerate}[label=(\arabic*)]
\item $A\in D_{\et}(X,\Lambda)$ 是 ULA 的。
\srcpage{101}
\item $A\in\mathrm{Fun}_{\cC_S}(X,S)$ 是左伴随。此时右伴随是 $\mathbb{D}_{X/S}(A)\in\mathrm{Fun}_{\cC_S}(S,X)$。
\item $p_1^*\mathbb{D}_{X/S}(A)\overset{L}{\otimes}_\Lambda p_2^*A\xrightarrow{\ \sim\ }R\mathrm{Hom}_\Lambda(p_1^*A,Rp_2^!A)$。
\end{enumerate}
\end{theorem}

\begin{proof}[证明概要]
对 (1) $\Rightarrow$ (2)，可以直接尝试构造单位与余单位。映射 $\beta$ 由定义同义反复地得到，但需要 (3) 来构造 $\alpha$。
\end{proof}

我们采用第二种进路，因为它很适合分析层范畴上的作用。
